\documentclass[11pt]{article}
\pdfinfoomitdate=1
\pdftrailerid{}
\pdfsuppressptexinfo=15
\usepackage[T1]{fontenc}
\usepackage{lmodern,microtype}
\usepackage[a4paper,margin=29mm]{geometry}
\usepackage{amsmath,amssymb,amsthm,mathtools}
\usepackage{enumitem,booktabs,array,longtable}
\usepackage[dvipsnames]{xcolor}
\usepackage[colorlinks=true,linkcolor=MidnightBlue,citecolor=MidnightBlue,urlcolor=MidnightBlue]{hyperref}
\usepackage[nameinlink,noabbrev]{cleveref}
\crefname{lemma}{lemma}{lemmas}
\Crefname{lemma}{Lemma}{Lemmas}
\crefname{proposition}{proposition}{propositions}
\Crefname{proposition}{Proposition}{Propositions}
\hypersetup{pdftitle={Sharp exceptional density for rowmotion on interval-closed sets of a rectangle},pdfauthor={Columbus Labs},pdfsubject={Fixed-width sharp exceptional-state asymptotics},pdfkeywords={rowmotion, interval-closed sets, noncrossing matchings, Lean}}
\usepackage{fancyhdr}
\pagestyle{fancy}\fancyhf{}\fancyhead[L]{\small Sharp exceptional density for interval-closed rowmotion}\fancyhead[R]{\small Columbus Labs}\fancyfoot[C]{\thepage}
\setlength{\headheight}{14pt}
\setlist{itemsep=2pt,topsep=5pt}
\allowdisplaybreaks[2]
\newtheorem{theorem}{Theorem}[section]
\newtheorem{lemma}[theorem]{Lemma}
\newtheorem{proposition}[theorem]{Proposition}
\newtheorem{corollary}[theorem]{Corollary}
\theoremstyle{definition}\newtheorem{definition}[theorem]{Definition}
\theoremstyle{remark}\newtheorem{remark}[theorem]{Remark}
\newcommand{\IC}{\mathcal{IC}}
\newcommand{\Row}{\operatorname{Row}}
\newcommand{\Min}{\operatorname{Min}}
\newcommand{\Inc}{\operatorname{Inc}}
\newcommand{\Ceil}{\operatorname{Ceil}}
\newcommand{\per}{\operatorname{per}}
\newcommand{\Cat}{\mathrm{Cat}}
\newcommand{\supp}{\operatorname{supp}}
\newcommand{\OO}{\mathrm{O}}
\newcommand{\II}{\mathrm{I}}
\newcommand{\TT}{\mathrm{T}}
\newcommand{\one}{\mathbf{1}}
\newcommand{\src}[1]{\texttt{\detokenize{#1}}}
\title{Sharp exceptional density for rowmotion\protect\\on interval-closed sets of a rectangle}
\author{Columbus Labs\\\small\href{mailto:info@columbus-labs.com}{info@columbus-labs.com}}
\date{12 October 2026}
\begin{document}
\maketitle
\begin{abstract}
For fixed $m\ge2$, let $A_m(n)$ count the interval-closed subsets of $[m]\times[n]$, and let $E_m(n)$ count those whose least period under reverse-toggle rowmotion is different from $n+m+1$. We prove
\[
\frac{E_m(n)}{A_m(n)}=\frac{m(m-1)}n+O_m(n^{-2}).
\]
The proof combines two disjoint rotation models for typical states with an explicit family of nonreturning states. The first model uses noncrossing matchings with at most one short support gap; the second uses sparse noncrossing partitions with one three-element block. Exact overlapping boundary-wave formulas establish both models for the actual toggle action. A telescoping boundary-time identity proves nonreturn on separated unit-gap valley chambers. Counting states, including the proper-period losses, then gives the sharp coefficient. In particular the result proves Conjecture 4.1 of Lafreni\`ere, Lewis, McNicholas, Striker, and Welch for every fixed width, with width one treated separately. The main fixed-width theorem also has an end-to-end Lean certificate; the proof below is mathematical rather than an appeal to that certificate.
\end{abstract}
\vspace{3pt}
\noindent\small\textcopyright\ 2026 Columbus Labs. Licensed under \href{https://creativecommons.org/licenses/by/4.0/}{CC BY 4.0}.\normalsize
\tableofcontents
\section{Introduction and statement}
An interval-closed subset of a finite poset is an order-convex subset: if it contains $x$ and $y$, it contains every $z$ with $x\le z\le y$. A toggle changes membership at one element when the resulting set is still interval-closed, and otherwise fixes the set. To fix the composition convention, let $(x_1,\ldots,x_{mn})$ be a linear extension from minimal to maximal elements and let $t_x$ be the toggle at $x$. We apply $t_{x_{mn}}$ first and $t_{x_1}$ last. Thus, with function composition evaluated right to left,
\[
 R=t_{x_1}\circ t_{x_2}\circ\cdots\circ t_{x_{mn}}.
\] We use the coordinatewise order on $[m]\times[n]$ and write $R$ for this permutation of $\IC([m]\times[n])$.

For a state $I$, let $\per(I)$ be its least positive return time. Set
\[
 A_m(n)=|\IC([m]\times[n])|,\qquad
 E_m(n)=|\{I\in\IC([m]\times[n]):\per(I)\ne n+m+1\}|.
\]
These are numbers of states, not numbers of orbits. Put
\[
 c_m=\frac{\Cat_m}{(2m)!}=\frac1{m!(m+1)!},\qquad
 \Cat_m=\frac1{m+1}\binom{2m}{m}.
\]

\begin{theorem}\label{thm:main}
For every fixed integer $m\ge2$, as $n$ tends to infinity through all positive integers,
\begin{align}
 E_m(n)&=c_m m(m-1)n^{2m-1}+O_m(n^{2m-2}),\label{eq:main-count}\\
 \frac{E_m(n)}{A_m(n)}&=\frac{m(m-1)}n+O_m(n^{-2}).\label{eq:main-ratio}
\end{align}
In particular, the fraction of states of least period $n+m+1$ tends to one.
\end{theorem}

The width is fixed throughout. Neither the implicit constants nor the onset of the asymptotic estimate is claimed uniform in a growing $m$. The theorem does not classify every exceptional orbit. It counts both states that fail to return after $n+m+1$ steps and states whose least period is a proper divisor of that number.

\subsection{Relation to prior work}
Elder, Lafreni\`ere, McNicholas, Striker, and Welch introduced the interval-closed toggle dynamics and established its global description \cite{ELMSW}. Their chain orbit classification, \cite[Theorem 3.2, arXiv v2]{ELMSW}, already implies the width-one density statement recalled in \cref{prop:width-one}. Lafreni\`ere, Lewis, McNicholas, Striker, and Welch simplified that description, classified width two, and proposed the fixed-width density-one statement in Conjecture 4.1 of \cite{LLMSW}. Their Section 4.1 already suggests a separated-endpoint family, Catalan endpoint orders, and comparison with all interval-closed sets. We retain this attribution: the density strategy and the width-two case are prior work.

The row-nonempty enumeration used below is also prior work: \cite[Theorem 2.7, arXiv v1]{ELLMSW} restates the Narayana-number formula from \cite[Theorem 4.4, arXiv v2]{ELMSW}. Noncrossing rotation models belong to the classical rowmotion literature; see \cite{SW}. The contribution developed here is the explicit passage from rotation to interval-closed-set rowmotion when two boundary waves overlap, together with a matching sharp lower bound from actual nonreturn dynamics. The coefficient $2$ at $m=2$ is consistent with, and already recoverable from, the existing width-two classification. No absolute novelty or publication-priority claim follows from the bounded literature comparison accompanying this project.

\subsection{Architecture of the proof}
The proof separates dynamics from enumeration. Set $N=n+m+1$ and $D=3m+8$. We construct injective equivariant maps into the actual state set from:
\begin{itemize}
\item $X$: noncrossing $m$-edge matchings on the $N$-circle with at most one support gap below $D$, excluding edges between ambient-adjacent vertices;
\item $Y$: noncrossing partitions on the $N$-circle with one triple and $m-2$ pairs, all occupied-support gaps at least $D$.
\end{itemize}
Their images are disjoint for all sufficiently large $n$. Almost every source state has least rotation period $N$, and the exceptions have order at most $n^{2m-2}$. Their combined cardinality is
\[
 |X|+|Y|=c_m\bigl(n^{2m}+mn^{2m-1}\bigr)+O_m(n^{2m-2}).
\]
The total state count has second coefficient $m^2$ instead of $m$. Thus these two models give the upper bound in \eqref{eq:main-count}.

For the lower bound, mark a valley $BA$ in a Dyck endpoint word and require the corresponding physical gap to be one. Keep every other endpoint gap and both margins above a fixed width-dependent floor. The resulting states cannot return at time $N$. There are $\frac{m-1}{2}\Cat_m$ marked valleys, and stars and bars gives their total leading coefficient $c_m m(m-1)$. The proof of nonreturn uses a finite sequence of rigorously controlled boundary events; it requires no full classification of their least periods.

The longer local calculations are in the appendices. Their complete candidate phase arrays and cutoff verifications make the dynamical argument independent of finite experiments or the formal certificate.

\section{Endpoint coordinates and the actual global rule}\label{sec:endpoints}
Call a state \emph{row-nonempty} if it meets each chain $\{i\}\times[n]$. Such a state is uniquely represented by intervals $[a_i,b_i]$ satisfying
\begin{equation}\label{eq:antitone}
 1\le a_m\le\cdots\le a_1\le n,\quad
 1\le b_m\le\cdots\le b_1\le n,\quad a_i\le b_i.
\end{equation}
For example, if $a_i<a_{i+1}$, the two comparable endpoints in rows $i$ and $i+1$ would force the omitted point $(i+1,a_i)$; this proves the first monotonicity condition. The upper condition follows in the same way. Conversely, if $(i,x)\le(k,z)\le(j,y)$ with both endpoints in the state, then $a_k\le a_i\le x\le z\le y\le b_j\le b_k$, proving convexity.

For a subset $S$ write $\Delta S$ and $\nabla S$ for its downward and upward closures, and $\Inc(S)$ for the elements incomparable with every member of $S$. The global rowmotion rule of \cite[Theorem 1.1]{LLMSW} is
\begin{equation}\label{eq:global}
 R(I)=\bigl(\Inc(C)\setminus I\bigr)
 \cup\bigl(\Delta C\setminus\Delta(\Min(I)\cap\Delta C)\bigr),
 \qquad C=\Min(\nabla I\setminus I).
\end{equation}
We use this previously proved theorem. The formal companion separately proves it from the literal toggles.

For any state, including a state with empty rows, define
\begin{align*}
 D_i&=\max\bigl(\{c:(r,c)\in C,\ r\ge i\}\cup\{0\}\bigr),\\
 U_i&=\min\bigl(\{c:(r,c)\in C,\ r\le i\}\cup\{n+1\}\bigr),\\
 E_i&=\max\bigl(\{c:(r,c)\in\Min(I)\cap\Delta C,\ r\ge i\}\cup\{0\}\bigr).
\end{align*}
The row form of \eqref{eq:global} is
\begin{equation}\label{eq:row-rule}
 R(I)_i=[E_i+1,D_i]\ \cup\ \bigl([D_i+1,U_i-1]\setminus I_i\bigr).
\end{equation}
An interval with reversed endpoints is empty. The formula follows because row $i$ of $\Delta C$ is $[1,D_i]$, row $i$ of $\nabla C$ is $[U_i,n]$, and row $i$ of the down-closure removed in \eqref{eq:global} is $[1,E_i]$.

For a row-nonempty state, its ceiling rows are exactly
\[
 (i,b_i+1)\quad\text{with }b_i<n\text{ and }(i=1\text{ or }b_i<b_{i-1}),
\]
and its minimum rows are exactly $(i,a_i)$ with $i=1$ or $a_i<a_{i-1}$. In particular, if both endpoint lists are strict and $b_1<n$, then all cutoffs are ordinary and
\begin{equation}\label{eq:translation}
 R(I)_i=[a_i+1,b_i+1].
\end{equation}
The departures from translation are caused by missing ceiling or minimum rows; their propagation is the boundary-wave mechanism used below.

For the endpoint-word encoding in this paragraph, assume additionally that all $2m$ endpoint coordinates are pairwise distinct. Read them in increasing order, labeling a lower endpoint by $A$ and an upper endpoint by $B$. The result is a Dyck word: every initial segment contains at least as many $A$'s as $B$'s. Conversely, a Dyck word on strictly ordered coordinates gives valid rank-paired intervals. A \emph{valley} means an adjacent $BA$ in this word. This extra distinctness assumption belongs to the word encoding; strictness of the two lists separately still allows cross-list coincidences in \eqref{eq:translation} and the ordinary passage below. We use one-based row indices but zero-based gap indices.

\section{A matching model with one short support gap}\label{sec:matching}
Fix $m\ge2$, $N=n+m+1$, and $D=3m+8$. Vertices of the circle are labeled $1,\ldots,N$ clockwise; $\rho$ adds one modulo $N$. A support gap is the positive clockwise distance between consecutive occupied vertices, including the wrap gap.

Let $X=X_{m,N,D}$ consist of noncrossing matchings with $m$ edges, at most one support gap less than $D$, and no edge joining cyclically adjacent vertices. The remaining vertices are unoccupied. This is a rotation-invariant family of actual matchings; no ordering of the edges is additional data.

At a linear cut, label the smaller endpoint of each edge $A$ and the larger endpoint $B$. Stack pairing recovers the matching from this Dyck word. Write its opening and closing coordinates in decreasing order as $p_1>\cdots>p_m$ and $q_1>\cdots>q_m$. Let
\[
 C_0=\{1,\ldots,m-1\}.
\]
A matching has a \emph{clear cut} if its support avoids $C_0$. At such a cut define
\begin{equation}\label{eq:decoder}
 F_0(M)_i=[p_i-(m-1),\ q_i-(m+1)].
\end{equation}

\begin{lemma}\label[lemma]{lem:direct}
The direct decoder is a well-defined injective map from clear-cut matchings in $X$ to row-nonempty states with strict endpoint lists.
\end{lemma}
\begin{proof}
Rank balance gives $p_i<q_i$. If $q_i-p_i=1$, these two consecutive occupied endpoints form an adjacent peak $AB$ and are stack-paired, contradicting the adjacent-edge exclusion. Thus $q_i-p_i\ge2$, exactly the interval nonemptiness condition. The smallest occupied coordinate is an opening and is at least $m$; the smallest closing is at least two larger than an opening. Similarly every opening is at least two below a closing. Together with $q_1\le N$, these observations put all decoded endpoints in $[1,n]$. Both lists are strict. Recovering $p_i=a_i+m-1$, $q_i=b_i+m+1$ and stack-pairing recovers the original matching.
\end{proof}

The adjacent-edge exclusion is rotation-invariant. It is not asserted that every adjacent matched edge makes every cut invalid: rotating such an edge to $N-1,N$ already makes the first decoded row invalid, which is enough to explain the exclusion.

\subsection{Local passage identities}
The following three local calculations are proved in \cref{app:local}. They are identities for the actual rowmotion action.
\begin{lemma}[Boundary passage identities]\label[lemma]{lem:passages}
Let $M\in X_{m,N,D}$ have an occupied vertex at $N$ and a clear support cut. In particular, $M$ has at most one support gap below $D$ and has no edge joining cyclically adjacent vertices. The following assertions retain both restrictions.
\begin{enumerate}[label=(\roman*)]
\item If both adjacent support gaps are at least $D$, then
\[
 R^m F_0(M)=F_0(\rho^m M).
\]
The identity permits the single allowed short support gap elsewhere in the support, subject to the adjacent-edge exclusion.
\item If the greatest occupied coordinates are $N-\delta,N$, with $1\le\delta<D$ the unique short support gap, then
\[
 R^{m+\delta}F_0(M)=F_0(\rho^{m+\delta}M).
\]
If these vertices are matched, $\delta\ge2$; otherwise they are both closing endpoints, with nested incident edges.
\item In case (ii), if $\delta\ge m$, the intermediate iterates with $m\le t\le\delta$ also agree with $F_0(\rho^tM)$.
\end{enumerate}
\end{lemma}
The third assertion matters: when the first point has finished crossing the corridor before the second arrives, there are clear cuts inside the two-point passage. Compatibility only at the terminal time would not define the desired global map.

For completeness, the local bounds are uniform: take $H=m+6$. All nonexceptional raw endpoint separations are at least $D-2\ge H$, and the left margin at a passage start is at least $D-m+1\ge H$. The second margin is $D-m+1$, not $D-2$. The third occupied point lies at most at $N-\delta-D$, so it does not reach $N$ before the two-point passage finishes.

\subsection{Gluing around the circle}
\begin{proposition}\label[proposition]{prop:matching-embedding}
There is an injective map $F_X:X\longrightarrow\IC([m]\times[n])$ satisfying
\[
 F_X(\rho M)=R(F_X(M)),
\]
and agreeing with $F_0$ at every clear cut. Consequently it preserves least periods.
\end{proposition}
\begin{proof}
Parameterize rotations by integer time modulo $N$. Each occupied vertex has a start time when it reaches $N$. With no short gap, place an ordinary interval of length $m$ at each start. If a short pair has successive start times $s,s+\delta$, replace its two intervals by the single interval $[s,s+\delta+m]$. Every other consecutive start gap is at least $D>m$, so the interiors of these intervals are disjoint, including at the cyclic seam. There is no third start in a two-point interval. The unique short pair, if present, specifies its interval canonically; no choice of orbit basepoint is needed.

Outside the interiors, use $F_0$. At phase $t$ inside an interval, use $R^t$ applied to $F_0$ at its initial clear cut. \Cref{lem:passages} identifies the terminal values and every intermediate clear value with $F_0$. For $\delta<m$ no intermediate cut is clear: the two forbidden-time intervals $[1,m-1]$ and $[\delta+1,\delta+m-1]$ cover all interior integer times. For $\delta\ge m$ the remaining clear times are exactly $m\le t\le\delta$, already covered by part (iii).

Inside a passage, equivariance is built into the definition. At its end it follows from the terminal identity; between passages it follows from \eqref{eq:translation}. The same argument crosses the cyclic seam. This defines $F_X$ as a function of the current matching, even when its rotation orbit has a stabilizer, because the event intervals and their direct initial states are determined by that matching.

To prove injectivity, suppose $F_X(M)=F_X(M')$. The union of their supports has at most $4m$ points. Each point forbids exactly $m-1$ rotation parameters from giving a clear cut. Thus if $N>4m(m-1)$ there is a simultaneous clear rotation. Equivariance and \cref{lem:direct} then imply $M=M'$. This size condition is automatic when $X$ is nonempty: its $2m$ positive support gaps include at least $2m-1$ gaps of size at least $D$, whence
\[
 N\ge(2m-1)D+1>4m(m-1).
\]
If $X$ is empty there is nothing to prove.

Finally, source return implies image return by equivariance; image return implies source return by injectivity. Hence their least positive periods are equal.
\end{proof}

\section{A disjoint model with one three-element block}\label{sec:partitions}
Let $Y=Y_{m,N,D}$ be the noncrossing partitions of an occupied subset of the circle with one block of size three, $m-2$ blocks of size two, and all $2m-1$ occupied-support gaps at least $D=3m+8$. Draw an arc between consecutive elements of each block in linear order. A triple $u<v<w$ contributes arcs $(u,v),(v,w)$. Read the shared vertex $v$ as $B$ before $A$. The $m$ arcs give a Dyck word with a marked valley whose two endpoint coordinates coincide. Conversely, merge the incident edges at a marked valley of a stack matching. This recovers exactly one triple and the remaining pairs; no crossing is introduced because the identified endpoints were consecutive. Thus this encoding is bijective.

At a clear cut, use the same direct formula \eqref{eq:decoder}. Each endpoint list is strict, and there is exactly one cross-list equality among the support coordinates. In physical endpoints it becomes a gap $a_p=b_q+2$, with $p<q$. Every other difference is at least $D-2$. Rank-paired support coordinates are distinct and separated, so all rows are nonempty and lie in the grid. The shared-coordinate reconstruction proves direct injectivity.

\begin{lemma}[The gap-two passage]\label[lemma]{lem:gap-two}
Suppose strict physical endpoint lists have $b_1=n$, exactly one small pair $a_k=b_{k+1}+2$, and all other endpoint differences and the left margin at least $H=m+6$. Assume $k$ is the first crossing, so $a_i<b_{i+1}$ for $i<k$. Set
\[
 r=\min\{i:k<i\le m,\quad i=m\text{ or }a_i>b_{i+1}\}.
\]
After $m$ rowmotion steps the endpoints are
\begin{align}
 A'_i&=\begin{cases}a_i+m&i<k,\\a_{i+1}+m&k\le i<m,\\1&i=m,\end{cases}\label{eq:gap2A}\\
 B'_i&=\begin{cases}b_{i+1}+m&i<r,\\a_r+m-2&i=r,\\b_i+m&i>r.\end{cases}\label{eq:gap2B}
\end{align}
At every interior time $j$, the only repeated-endpoint position is between rows $j,j+1$; the upper list, both lists, or the lower list repeats there, in that order.
\end{lemma}
\begin{proof}
For $0\le j<k$ use the ordinary upper wave: $A_i=a_i+j$ and $B_i=b_{i+1}+j$ for $i\le j$, otherwise $b_i+j$. At $j=k-1$ the missing ceiling row is $k$. Its cutoffs are
\[
 D_k=b_{k+1}+j+1=a_k+j-1,\qquad E_k=a_{k+1}+j.
\]
The incomparable contribution is removed by the old interval, leaving the shifted next row. For $k\le j<r$, the resulting arrays are
\[
 A_i=\begin{cases}a_{i+1}+j&k\le i\le j,\\a_i+j&\text{otherwise},\end{cases}
 \quad
 B_i=\begin{cases}b_{i+1}+j&i\le j,\\b_i+j&i>j.\end{cases}
\]
Both lists repeat at $j,j+1$. In row $h=j+1$ the cutoffs are $D_h=b_{h+1}+j+1$ and $E_h=a_{h+1}+j$. Until $h=r$, the overlap $a_h<b_{h+1}$ makes the incomparable term empty, so the equal-row pair propagates. At $h=r<m$, its large crossing instead gives the interval $[a_{r+1}+j+1,a_r+j-1]$. The upper list then has the final sources in \eqref{eq:gap2B}; the lower repeated endpoint propagates to the bottom, which inserts $1$. If $r=m$, the zero cutoffs give $[1,a_m+m-2]$ directly.

All changed row sources are original valid rows or the displayed crossing interval. Separation makes the crossing interval nonempty, prevents unintended equalities, and keeps all retained upper endpoints below $n$. The exceptional $a_k$ is deleted before it can create an additional comparison. This proves the arrays, their validity, and the terminal formulas by \eqref{eq:row-rule}.
\end{proof}

\begin{proposition}\label[proposition]{prop:partition-embedding}
There is an injective equivariant map $F_Y:Y\longrightarrow\IC([m]\times[n])$, agreeing with its direct decoder at every clear cut and preserving least periods.
\end{proposition}
\begin{proof}
At a start with an occupied vertex at $N$, only that vertex crosses the cut within $m$ steps. If its block is a pair, the ordinary wave of \cref{app:ordinary} applies with the shared vertex passive: the flipped opening is not the marked opening, so the small gap is not at its switch. If its block is a triple, write the expanded word as
\[
 U_0 A W\,B A\,V B,
\]
where the displayed $BA$ is shared and $U_0,W,V$ are Dyck words. Its last primitive opening gives the first crossing $k$; the mate of the preceding marked $B$ gives the next crossing $r$. Under rerooting the triple, the shared coordinate moves from its middle vertex to its first vertex. The new word is $A U_0 B A W B V$, with the displayed $B|A$ marked. The coordinate changes are exactly \eqref{eq:gap2A}--\eqref{eq:gap2B}. Thus $R^mF_0(P)=F_0(\rho^mP)$ in both cases.

At most one occupied point lies in the corridor, since $D>m$. If it is $j\in C_0$, define $F_Y(P)=R^jF_0(\rho^{-j}P)$; otherwise use $F_0(P)$. The terminal identity and translation prove equivariance. The common-clear-cut proof used for $X$ also gives injectivity: two supports forbid at most $2(2m-1)(m-1)$ shifts, whereas nonemptiness of $Y$ implies $N\ge(2m-1)D$ and hence leaves a clear shift. At that shift the direct decoder is injective. Least periods are preserved as before.
\end{proof}

\begin{lemma}\label[lemma]{lem:disjoint}
If $N>(4m-1)(m-1)$, the images of $F_X$ and $F_Y$ are disjoint.
\end{lemma}
\begin{proof}
If two images coincide, equivariance preserves their equality after any common rotation. Their support union has at most $4m-1$ points, so some rotation makes both cuts clear. Recovering $p_i=a_i+m-1$, $q_i=b_i+m+1$ from the common state would then give both no cross-list equality (a matching) and exactly one cross-list equality (a triple partition), a contradiction.
\end{proof}

\section{A codimension-one family of nonreturning states}\label{sec:nonreturn}
The upper bound needs many states returning at the baseline time; the lower bound needs states that definitely do not. Here a single physical unit gap at a valley supplies the obstruction.

Let $d=2m-1$ and set
\begin{equation}\label{eq:floor}
 G=G(m)=m+6+(m+2)d(d+1).
\end{equation}
In a strict endpoint word mark one valley $BA$ and give its two coordinates distance one. Require every other consecutive endpoint gap, the one-based left margin, and the right margin to be at least $G$. The Dyck condition produces an actual row-nonempty state. The word, mark, and gaps are uniquely recovered from the state, since there is exactly one unit gap.

\subsection{Exact strict-boundary events}
A \emph{strict boundary state} has strict endpoint lists and $b_1=n$. For the unit-gap events considered here, assume also that all $2m$ coordinates are distinct. List them as $z_1<\cdots<z_{2m}=n$, and mark the valley pair $z_{j_*}=b_q$, $z_{j_*+1}=a_p=b_q+1$. Define its separation parameter by
\[
 \eta(I)=\min\Bigl(\{z_1\}\cup
 \{z_{j+1}-z_j:1\le j<2m,\ j\ne j_*\}\Bigr).
\]
Thus the left margin is the one-based coordinate $z_1=a_m$, and the zero right margin is omitted. The condition $\eta(I)\ge H$ says precisely that every unmarked consecutive gap and this left margin are at least $H$; it does not require a gap between the marked pair. Suppose $H\ge m+6$ and $\eta(I)\ge H$, and put
\[
 k=\min\{i:i=m\text{ or }a_i>b_{i+1}\}.
\]
There are three cases:
\begin{itemize}
\item $\OO$ (passive): $(p,q)\ne(k,k+1)$;
\item $\II$ (interior): $(p,q)=(k,k+1)$ and $k\ge2$;
\item $\TT$ (top): $(p,q)=(k,k+1)=(1,2)$.
\end{itemize}
In case $\OO$, $k\ne p$: if $q=p+1$, first crossing gives $k\le p$ and passivity excludes equality; if $q>p+1$, the separated endpoint $b_{p+1}$ is larger than $a_p$. Thus the short pair is not removed by the ordinary wave.

\begin{lemma}[Unit-gap boundary events]\label[lemma]{lem:events}
Under the preceding hypotheses, the local durations are $\ell=m,m+1,m+3$ in cases $\OO,\II,\TT$, respectively. Their terminal strict endpoint lists are:
\begin{align*}
\OO:\quad A'&=(a_1+m,\ldots,a_{k-1}+m,a_{k+1}+m,\ldots,a_m+m,1),\\
B'&=(b_2+m,\ldots,b_k+m,a_k+m-2,b_{k+1}+m,\ldots,b_m+m);\\[2pt]
\II:\quad A'&=(a_1+m+1,\ldots,a_{k-1}+m+1,a_{k+1}+m+1,\ldots,a_m+m+1,1),\\
B'&=(b_2+m+1,\ldots,b_k+m+1,a_{k-1}+m,b_{k+1}+m+1,\ldots,b_m+m+1);\\[2pt]
\TT:\quad A'&=(a_2+m+3,\ldots,a_m+m+3,1),\\
B'&=(b_2+m+3,\ldots,b_m+m+3,a_m+m+2).
\end{align*}
Their top upper endpoints are below $n$. Translation to the next strict boundary gives total elapsed time
\begin{equation}\label{eq:event-time}
 \tau=\begin{cases}n-a_1+2&\OO\text{ with }k=1,\\n-b_2&\text{otherwise}.
 \end{cases}
\end{equation}
There is no intervening strict boundary state. At the next strict boundary state $I^+$, with the transported marked pair described below, $\eta(I^+)\ge H-(m+2)$.
\end{lemma}
\begin{proof}
The passive case is the ordinary wave in \cref{app:ordinary}, because the only small gap is passive and $k\ne p$. The other two cases, including their empty-row phases, are proved directly in \cref{app:unit}. These formulas give \eqref{eq:event-time} on adding $n-B'_1$ translation steps. During each local passage there is a repeated endpoint or an empty row. The top passage revisits $b_1=n$ only with a repeated lower endpoint. Thus no earlier strict boundary state occurs.

All nonmarked endpoint differences at the terminal phase are at least $H-2$ (and at least $H-1$ in the collision cases). The inserted $1$ becomes $1+n-B'_1$ at the next boundary. The displayed upper sources give this new left margin at least $H-\ell+1\ge H-m-2$. This proves the conservative common loss bound.
\end{proof}

Write the marked Dyck word uniquely as $UAVB$, with $AVB$ its last primitive component. The last-component opening is the first crossing $a_k$. Reading the endpoint formulas gives the following marked-word update:
\begin{align*}
\OO:&\quad UAVB\longmapsto AUBV,\quad\text{carry the old marked letters};\\
\II:&\quad UAVB\longmapsto AUBV,\quad\text{mark }B\mid\text{the first }A\text{ of }V;\\
\TT:&\quad UAB\longmapsto ABU,\quad\text{mark the displayed }B\mid\text{the first }A\text{ of }U.
\end{align*}
In the latter two cases the old mark is at $U|A$, so $U$ is nonempty, and in the interior case $V$ is nonempty. These descriptions follow from which opening is deleted or flipped; they do not assume a return time. The zero-based position of the marked $B$ changes by $1,2,-(d-3)$ in the three cases.

\subsection{Gap transport and a closed-return identity}
At a strict boundary state list the $2m$ endpoint coordinates increasingly. Delete the marked unit gap from the list of gaps, keep the one-based left gap, and omit the zero right margin. This gives $h=(h_0,\ldots,h_{d-1})$ with
\[
 S=\sum_i h_i=n-1.
\]
Let $\sigma(h)=(h_{d-1},h_0,\ldots,h_{d-2})$, and let $e_i$ denote coordinate unit vectors.

\begin{lemma}\label[lemma]{lem:gap-transport}
At successive strict boundary states, the gaps obey
\begin{align}
\OO:\quad h'&=\sigma(h-2e_s+2e_{s+1})-(m-1)e_0+(m-1)e_1,\label{eq:O-gap}\\
\II:\quad h'&=\sigma h-me_0+me_1,\label{eq:I-gap}\\
\TT:\quad h'&=\sigma h-(m+1)e_0+(m+1)e_1.\label{eq:T-gap}
\end{align}
Here $s\in\{0,\ldots,d-2\}$ in case $\OO$ is the large gap immediately preceding the flipped opening. The elapsed time is $\tau=h_{d-1}+\epsilon$, where $\epsilon=2$ for $\OO$ with $s=d-2$, $\epsilon=1$ for $\TT$, and $\epsilon=0$ otherwise.
\end{lemma}
\begin{proof}
In the passive case, first move the flipped opening left by two: its two neighboring large gaps change by $-2,+2$. Neither is the marked gap because the mark is passive. Removing the final endpoint, inserting the new first endpoint $1$, and translating to the next boundary rotates these gaps, with an extra transfer of $m-1$ from the new first gap to the next. This is \eqref{eq:O-gap}. Only when the flipped opening is the penultimate endpoint, equivalently $s=d-2$ or $k=1$, does moving it change the last surviving endpoint; the event time gains two.

In case $\II$, the marked opening is deleted and a new closing is inserted one unit before the next opening. Removing the old unit gap and inserting the new one cancel on the intervening large gap. The other large internal gaps retain their order. Since $\ell=m+1$, the new first and second gaps are $h_{d-1}-m$ and $h_0+m$, giving \eqref{eq:I-gap}. In case $\TT$, the final opening and closing are deleted and a closing is inserted one unit before the smallest surviving opening. The same difference computation, with $\ell=m+3$ and time $n-b_2=h_{d-1}+1$, gives \eqref{eq:T-gap}. All endpoint orders used in these computations are guaranteed by \cref{lem:events}.
\end{proof}

Put $Q(h)=\sum_{i=0}^{d-1}ih_i$. Since
\[
 Q(\sigma h)-Q(h)=S-dh_{d-1},
\]
\cref{lem:gap-transport} gives the identity
\begin{equation}\label{eq:cocycle}
 d\tau=S+\kappa-\bigl(Q(h')-Q(h)\bigr),\qquad
 \kappa_{\OO}=m+1,\quad\kappa_{\II}=m,\quad\kappa_{\TT}=3m.
\end{equation}
For $\OO$, the two-unit transfer contributes $2$ to $Q$ after the cyclic shift, except at $s=d-2$, where it contributes $-2(d-1)$; the additional $d\epsilon=2d$ cancels that discrepancy. The other two values follow directly.

Suppose a closed segment has $K>0$ events, including $O$ passive, $I$ interior, and $T$ top events. Equality of the initial and final marked positions yields
\[
 O+2I=(d-3)T,\qquad K+I=(d-2)T.
\]
In particular $T>0$, since $T=0$ forces $O=I=0$. Summing \eqref{eq:cocycle}, with $N=S+m+2$, gives the key identity for the actual elapsed time $t$:
\begin{equation}\label{eq:closed-return}
 dt=KN+2T.
\end{equation}
No coprimality of a marked-word cycle is needed.

\begin{proposition}\label[proposition]{prop:nonreturn}
Every separated marked unit-gap chamber with free gaps at least $G(m)$ consists of states $I$ satisfying $R^N(I)\ne I$.
\end{proposition}
\begin{proof}
First suppose the initial state is strict boundary and all its $d$ large gaps are at least $G$. \Cref{lem:events} applies to the first $d+1$ events, because after $j\le d$ events the separation is at least $G-j(m+2)\ge m+6$.

The affine gap formulas have the form $h^{j+1}=\sigma h^j+v_j$, every coordinate of $v_j$ being at least $-(m+2)$. Hence
\[
 h_i^j\ge h^0_{i-j\bmod d}-j(m+2).
\]
As each event time is at least the current last gap, the first $d+1$ event times sum to at least
\[
 S+h^0_{d-1}-\frac{(m+2)d(d+1)}2
 \ \ge\ S+G-\frac{(m+2)d(d+1)}2\ >\ S+m+2=N.
\]
Here the last-gap coordinate has consumed one full cycle of initial gaps and one repeated gap. Thus, if the initial state returned at time $N$, it would do so at one of the first $K\le d$ strict-boundary events. Equation \eqref{eq:closed-return} would imply
\[
 (d-K)N=2T.
\]
For $K=d$ this contradicts $T>0$. For $K<d$ it implies $N\le2T\le2K\le2(d-1)$, while $N\ge dG+m+2$, again impossible.

For an interior chamber, translate until $b_1=n$. Its two outer margins merge into the boundary left gap; every other large gap is unchanged. The boundary floor still holds. Since $R$ is a permutation and its iterates commute, a baseline return of the original state would imply a baseline return of this boundary state. The latter has just been excluded.
\end{proof}

\paragraph{The role of the floor.}
The value $G(m)$ is a conservative proof convenience. It keeps the local event hypotheses valid through the bounded sequence used above and makes the elapsed-time estimate exceed $N$. No optimal separation threshold is claimed. The proof does not extend the nonreturn assertion to zero or arbitrarily small unmarked gaps; finite tests in such regimes would not establish that extension.

\section{Counting the models and completing the asymptotic}\label{sec:counting}
All estimates in this section hold at fixed $m$ and fixed spacing constants. Let
\[
 V_m=\sum_{w\text{ Dyck of semilength }m}\#\{\text{valleys of }w\}
     =\frac{m-1}{2}\Cat_m.
\]
For example, the Narayana distribution of the number of peaks is symmetric under $k\leftrightarrow m+1-k$; its mean is $(m+1)/2$. Valleys equal peaks minus one, giving the displayed identity. This is standard Dyck enumeration.

\subsection{All actual states}
The known row-nonempty count is
\begin{equation}\label{eq:narayana}
 F_m(n)=\frac1{m+1}\binom{n+m-1}{m}\binom{n+m}{m}
       =c_m(n)_m(n+1)_m,
\end{equation}
Here $(x)_m=x(x+1)\cdots(x+m-1)$ is the \emph{rising factorial}, with $(x)_0=1$. The formula is \cite[Theorem 2.7, arXiv v1]{ELLMSW}, which attributes it to \cite[Theorem 4.4, arXiv v2]{ELMSW}. A state with an empty row can be overcounted by choosing that row and an arbitrary interval or empty set in each other row. Thus their number is at most
\[
 m\left(1+\frac{n(n+1)}2\right)^{m-1}=O_m(n^{2m-2}).
\]
The two rising factorials have sum of shifts $m(m-1)/2+m(m+1)/2=m^2$, so
\begin{equation}\label{eq:denominator}
 A_m(n)=c_m\bigl(n^{2m}+m^2n^{2m-1}\bigr)+O_m(n^{2m-2}).
\end{equation}
This includes all empty-row states; none are silently removed from the denominator.

\subsection{Matching states}
There are $\Cat_m\binom N{2m}$ unrestricted noncrossing $m$-edge matchings. Supports with at least two gaps below $D$ number $O_m(N^{2m-2})$. Indeed, two disjoint close pairs are specified by two starting vertices, bounded offsets, and $2m-4$ other vertices. If the pairs share a vertex, specify that vertex, two bounded offsets, and $2m-3$ others. There are only width-dependent numbers of choices of positions and noncrossing pairings. This also counts gaps crossing the cyclic seam.

The number of matchings with a \emph{marked} ambient-adjacent edge is exactly
\[
 N\Cat_{m-1}\binom{N-2}{2m-2}.
\]
Choose its cyclic first vertex, choose all other support vertices, and match them noncrossingly; an adjacent chord crosses none of the others. Multiple marked adjacent edges occur only on the preceding two-short-gap locus. Therefore marking versus not marking changes only the error term, and
\begin{align}
 |X|&=\Cat_m\binom N{2m}-N\Cat_{m-1}\binom{N-2}{2m-2}
       +O_m(n^{2m-2})\notag\\
 &=c_m\bigl(n^{2m}+(2m-m^2)n^{2m-1}\bigr)+O_m(n^{2m-2}).\label{eq:X-count}
\end{align}
The coefficient calculation uses $N=n+m+1$,
\[
 \Cat_m\binom N{2m}=c_m(n^{2m}+3mn^{2m-1})+O_m(n^{2m-2}),
 \qquad \frac{\Cat_{m-1}}{(2m-2)!}=m(m+1)c_m.
\]

\subsection{Triple partitions}
The marked-valley bijection of \cref{sec:partitions} gives $V_m$ partitions on each $(2m-1)$-element support. Losing supports with a short gap costs $O_m(N^{2m-2})$ by specifying one close pair. Hence
\begin{equation}\label{eq:Y-count}
 |Y|=V_m\binom N{2m-1}+O_m(n^{2m-2})
     =c_m m(m-1)n^{2m-1}+O_m(n^{2m-2}).
\end{equation}
The supports and their partitions are counted once; no extra phase or rooted-word multiplicity occurs.

\subsection{Proper-period losses are state losses}
\begin{lemma}\label[lemma]{lem:proper}
The number of source states in $X$ with least rotation period less than $N$ is $O_m(N^m)$. The analogous number in $Y$ is $O_m(N^{\lfloor(2m-1)/2\rfloor})$. Both are $O_m(n^{2m-2})$ for $m\ge2$.
\end{lemma}
\begin{proof}
If a support of size $k$ has a nontrivial stabilizer, it is invariant under a cyclic subgroup of some order $q>1$ dividing both $N$ and $k$. The action on the circle is free, so the support is the union of $k/q$ orbits of that subgroup. There are $\binom{N/q}{k/q}$ choices. For fixed $k$ there are at most as many possible $q$ as divisors of $k$, independent of $N$. The number of matchings or partitions on a fixed support is bounded by a constant depending only on $k$. Consequently the total is $O_k(N^{\lfloor k/2\rfloor})$. Apply this with $k=2m$ and $k=2m-1$. The odd case could be sharpened using $q\ge3$, but that is unnecessary.
\end{proof}
In particular the matching loss is $O(n^2)$ at $m=2$, exactly the allowed error. The lemma counts actual source states with a proper return. It does not count orbits and then multiply by a possible orbit length.

\subsection{The upper and lower bounds}
Injectivity, equivariance, and disjointness of the two models, together with \cref{lem:proper}, give for all sufficiently large $n$
\begin{align*}
 A_m(n)-E_m(n)
 &\ge |X|+|Y|-O_m(n^{2m-2})\\
 &=c_m\bigl(n^{2m}+mn^{2m-1}\bigr)+O_m(n^{2m-2}).
\end{align*}
Subtracting from \eqref{eq:denominator} proves
\[
 E_m(n)\le c_m m(m-1)n^{2m-1}+O_m(n^{2m-2}).
\]

For the lower bound, each marked unit-gap chamber has $2m$ free gaps: the left margin, the right margin, and the $2m-2$ unmarked internal gaps. They sum to $n-1$. With each at least $G$, their exact number, for $n\ge1+2mG$, is
\[
 \binom{n-2mG+2m-2}{2m-1}.
\]
The $V_m$ choices of word and marked valley give disjoint actual states. By \cref{prop:nonreturn}, every one is exceptional. Thus
\begin{align*}
 E_m(n)&\ge V_m\binom{n-2mG+2m-2}{2m-1}\\
 &=c_m m(m-1)n^{2m-1}+O_m(n^{2m-2}).
\end{align*}
This is a full polynomial-error lower estimate, not merely a liminf. The two bounds prove \eqref{eq:main-count}. Dividing by \eqref{eq:denominator}, whose positive leading coefficient is $c_m$, proves \eqref{eq:main-ratio} and completes \cref{thm:main}.

\subsection{Width one and the original positive-width conjecture}
\begin{proposition}[Consequence of the chain classification in {\cite[Theorem 3.2, arXiv v2]{ELMSW}}]\label[proposition]{prop:width-one}
The fraction of interval-closed subsets of $[1]\times[n]$ of least rowmotion period $n+2$ tends to one.
\end{proposition}
\begin{proof}
The cited theorem already gives this limit. For completeness, the following two-point model rederives the needed bound. Put $N=n+2$ and label the circle $1,\ldots,N$. Non-cyclic-adjacent unordered pairs $\{p<q\}$ correspond to intervals $[p,q-2]$. They give every nonempty interval except $[1,n]$. Rotation by one agrees with rowmotion, since
\[
 R([a,b])=[a+1,b+1]\ (b<n),\qquad
 R([a,n])=[1,a-1]\ (a\ge2).
\]
These formulas follow directly from \eqref{eq:global} on a chain. The two omitted states, empty and full, interchange. A two-point subset has a nontrivial rotation stabilizer only when its points are antipodal. Therefore at most $N/2+2$ of the $1+n(n+1)/2$ states can have least period different from $N$, which proves the limit.
\end{proof}
\Cref{thm:main}, together with the previously known width-one case recalled in \cref{prop:width-one}, establishes Conjecture 4.1 of \cite{LLMSW} for every fixed positive width. The sharp coefficient formula is stated only for $m\ge2$; antipodal symmetries explain why substituting $m=1$ into that formula is not justified.

\section{Formal certificate, reproducibility, and limits}\label{sec:certificate}
The main result is accompanied by the Lean declaration
\begin{center}
\src{RowmotionSharpAllWidth.actualRowmotion_sharp_all_width}.
\end{center}
Its only public hypotheses are a natural number $m$ and $2\le m$. The definitions count literal order-convex predicate subsets of $\mathrm{Fin}\,m\times\mathrm{Fin}\,n$, and states lacking least period $n+m+1$ under reverse row-major toggles. The row-major list is proved to be a duplicate-free complete linear extension. There is no remaining count, injectivity, equivariance, disjointness, or nonreturn assumption in the final theorem.

The \href{https://github.com/ColumbusLabs/interval-closed-rowmotion}{companion repository} supplies sources in \src{formal/sharp-allwidth}. Two inventories bind their contents:
\begin{itemize}
\item \src{SOURCE_MANIFEST.json} lists the SHA-256 hashes of all 225 mathematical source modules.
\item \src{PACKET_MANIFEST.json} additionally binds the verification drivers, build configuration, and source-checking scripts.
\end{itemize}
The declaration above is the final unconditional theorem. The repository's theorem-to-source map aligns the written proof with the corresponding Lean declarations.

The recorded verification rebuilt all 225 local mathematical modules, audited all 9,278 project declarations (including generated declarations), and replayed the complete imported environment into a fresh pinned Lean kernel. Its permitted axiom dependencies are the standard \src{propext}, \src{Classical.choice}, and \src{Quot.sound}; no project axiom or admitted proof was accepted. It used Lean 4.34.1, commit \src{5045d0056413266e57c625dcd7c365b10e377c52}, and mathlib revision \src{d13f23b723b8a846827a245b89c10fc7d3f11612}. The replay is a fresh run of the pinned kernel, not an independently implemented proof checker. The replay and audit drivers are separate from the mathematical proof imports.

The formal denominator proof uses an explicit endpoint/Dyck equivalence rather than assuming \eqref{eq:narayana}: weak coordinates with a stable tie convention give the exact sum
\[
 \sum_{w\text{ Dyck of semilength }m}
 \binom{n+2m-1-\operatorname{val}(w)}{2m}.
\]
It then derives the same two leading terms and handles empty rows separately. This alternative formal route does not turn the previously known enumeration into a novelty claim. Likewise the formal proper-period argument uses a finite support-permutation encoding, while \cref{lem:proper} uses the elementary cyclic-subgroup count; both bound actual source states to the required order.

The accompanying theorem-to-source map records these alignments and the precise local modules for the appendices. Build instructions, source hashes, toolchain pins, and a verification summary accompany the manuscript. Recompiling this PDF does not rerun the Lean certificate; those are distinct reproducibility tasks.

\paragraph{Scope left open.}
The proof establishes a sharp state-density asymptotic at each fixed width. It does not give a uniform growing-width error bound, a complete orbit classification, or general least periods for noncoprime boundary-phase cycles. The separately certified width-three period-distribution refinement is outside this manuscript's proof scope. Questions about lower-order coefficients and growing-width sharpness require additional analysis.

\appendix
\section{Boundary waves for the rotation embeddings}\label{app:local}
This appendix proves \cref{lem:passages}. The candidate arrays below are explicit formulas checked by the actual cutoff rule \eqref{eq:row-rule}; composing their array substitutions is not an assumption that overlapping rowmotion waves commute.

\subsection{An isolated passage with a passive close pair}\label{app:ordinary}
Suppose $a_i,b_i$ are strict decreasing valid endpoint lists with $b_1=n$, and let $k$ satisfy
\[
 a_i\le b_{i+1}\quad(i<k),\qquad
 a_k\ge b_{k+1}+3\quad(k<m).
\]
Assume $a_m\ge1$, $b_2+m<n$ (when $m\ge2$), and $a_k+m-2<n$. For $0\le t<k$, define
\begin{equation}\label{eq:ordinary-upper}
 A_i(t)=a_i+t,\qquad
 B_i(t)=\begin{cases}b_{i+1}+t&i\le t,\\b_i+t&i>t.\end{cases}
\end{equation}
For $k\le t<m$, define
\begin{align}
 A_i(t)&=\begin{cases}a_{i+1}+t&k\le i\le t,\\a_i+t&\text{otherwise},\end{cases}\label{eq:ordinary-lower}\\
 B_i(t)&=\begin{cases}b_{i+1}+t&i<k,\\a_k+t-2&i=k,\\b_i+t&i>k.\end{cases}\notag
\end{align}
The terminal state is
\begin{align}
 A_i(m)&=\begin{cases}a_i+m&i<k,\\a_{i+1}+m&k\le i<m,\\1&i=m,\end{cases}\label{eq:ordinary-terminal}\\
 B_i(m)&=\begin{cases}b_{i+1}+m&i<k,\\a_k+m-2&i=k,\\b_i+m&i>k.\end{cases}\notag
\end{align}

Here is the full cutoff induction. Before the switch, only ceiling row $h=t+1$ is missing. For $h<k$, its next ceiling gives $D_h=b_{h+1}+t+1>A_h$, and its minimum qualifies. Row $h$ becomes $[a_h+t+1,b_{h+1}+t+1]$, while all other rows translate. At $h=k<m$, $A_h\ge D_h+2$, so that minimum fails to qualify, and $E_h=a_{h+1}+t$. The incomparable left interval joins the part below the ceiling, giving $[a_{k+1}+t+1,a_k+t-1]$. At $k=m$, the zero cutoffs give $[1,a_m+m-2]$. After switching, every ceiling row is present and only minimum row $t+1$ is missing. Its $E$ cutoff reads the next lower endpoint, or $0$ at the bottom. This propagates the lower equality and gives \eqref{eq:ordinary-terminal}.

All rows are valid: before switching use $a_i\le b_{i+1}$; at the switch use $a_{k+1}\le b_{k+1}\le a_k-3$; afterwards each lower source moves downward. The flipped upper endpoint is strictly between its neighbors because $b_k\ge a_k>a_k-2>b_{k+1}$. All remaining list comparisons are inherited. The strict top bounds exclude new ceiling defects. Thus the formulas allow singleton rows and arbitrarily small positive same-list gaps.

For an isolated matching passage, the last closing $N$ is matched to the first balanced-suffix opening $p_k$. The stack rule gives $p_i<q_{i+1}$ for $i<k$, and $p_k>q_{k+1}$ when $k<m$. An opening and closing at distance one would be an excluded adjacent matched peak, so decoding yields $a_i\le b_{i+1}$ before the switch and $a_k\ge b_{k+1}+3$ at it. Large neighboring support gaps give the strict top bounds and $a_m\ge D-m+1$. Under $\rho^m$, $N$ becomes the new opening $m$, while $p_k$ changes from opening to closing. The two decoder offsets differ by two, so the resulting endpoints are precisely \eqref{eq:ordinary-terminal}. This proves \cref{lem:passages}(i), including every passive location of the close pair.

\subsection{A short matched leaf}\label{app:leaf}
Suppose $m\ge2$, $\delta\ge2$, $H\ge m+6$, and
\[
 b_1=n,\qquad a_1=n-\delta+2,\qquad a_m\ge H.
\]
All other endpoints are at least $H$ below $a_1$, and pairwise separated by at least $H$. Set $T=m+\delta$. At time $0$ use the input. For $1\le t<T$, define
\begin{equation}\label{eq:leaf-A}
 A_i(t)=\begin{cases}
 a_{i+1}+t&i\le t<m,\\
 a_i+t&t<m,\ i>t,\\
 a_{i+1}+t&t\ge m,\ i<m,\\
 t-m+1&t\ge m,\ i=m.
 \end{cases}
\end{equation}
Put $u=\max(0,t-\delta)$, and define
\begin{equation}\label{eq:leaf-B}
 B_i(t)=\begin{cases}
 b_{i+1}+t&i\le u,\\
 n-\delta+t&u=0,\ i=1,\\
 b_i+t&\text{otherwise}.
 \end{cases}
\end{equation}
Then the terminal state is
\begin{equation}\label{eq:leaf-terminal}
 (A_i(T),B_i(T))=
 \begin{cases}(a_{i+1}+T,b_{i+1}+T)&i<m,\\(1,\delta-1)&i=m.
 \end{cases}
\end{equation}

At time $0$ the only missing ceiling is row $1$, with $D_1=b_2+1<a_1$ and $E_1=a_2$. Its output is $[a_2+1,a_1-1]$, giving the first lower equality and \eqref{eq:leaf-B}. For $1\le t<m$, the missing minimum is row $l=t+1$. Its next qualifying minimum has source $a_{l+1}+t$, or the zero sentinel if $l=m$. Its ceiling is present, so the lower equality advances as in \eqref{eq:leaf-A}.

Before $t=\delta$, all ceilings are present. At $t=\delta$ the first upper endpoint reaches $n$; thereafter the missing ceiling is row $h=u+1$. Until the last step, $h<m$ and
\[
 A_h(t)=a_{h+1}+t,\qquad D_h=b_{h+1}+t+1,\qquad E_h=A_h(t).
\]
The possible missing minimum lies $\delta\ge2$ rows lower, so the displayed minimum is present. As $a_{h+1}\le b_{h+1}$, the incomparable contribution is removed. Thus the upper equality advances. At the final input $t=T-1$, row $m$ has $A_m=\delta$ and missing ceiling, with $D_m=E_m=0$. Its output is the complement of $[\delta,B_m]$ in $[1,B_m]$, namely $[1,\delta-1]$.

These are all exceptional cutoffs. Every changed lower source comes from a later original row, or is the inserted $t-m+1$. The top-row flip is valid since $a_1-a_2\ge H$. Retained upper endpoints satisfy $b_2+T\le n+m+2-H<n$, and the special $n-\delta+t$ reaches $n$ only at its intended second passage. This proves validity and the full induction.

If the close pair in the matching is matched, it is the last $AB$ leaf. After rotating by $T$, its endpoints are $m,m+\delta$, so its decoded row is $[1,\delta-1]$, while all other endpoints translate. This is \eqref{eq:leaf-terminal}. For $\delta\ge m$, the formulas on $m\le t\le\delta$ give
\[
 A_i=a_{i+1}+t\ (i<m),\quad A_m=t-m+1,\quad
 B_1=n-\delta+t,\quad B_i=b_i+t\ (i>1),
\]
which is exactly the intermediate direct decoder. This proves parts (ii)--(iii) for the matched case.

\subsection{Two close closings from distinct edges}\label{app:distinct}
Now let $m\ge2$, $\delta\ge1$, and $H\ge m+6$. Start from valid strictly decreasing endpoint lists $a_i,b_i$ satisfying
\[
 b_1=n,\qquad b_2=n-\delta,\qquad a_m\ge H.
\]
Every other endpoint lies at most at $b_2-H$, and every pair other than $b_1,b_2$ is separated by at least $H$. Let
\[
 k=\min\{i:i=m\text{ or }a_i>b_{i+1}\},\qquad
 r=\min\{i:1\le i<k,\ a_i>b_{i+2}\},
\]
where $b_{m+1}=0$ is used only in defining $r$. Then $2\le k\le m$ and $1\le r<k$: the last eligible $i=k-1$ always satisfies the inequality. Put $T=m+\delta$.

\paragraph{Candidate phases.}
Subtract the common translation $t$ and first define older-wave arrays:
\begin{equation}\label{eq:older-x}
 x_i(t)=\begin{cases}
 a_i&t<k,\\
 a_{i+1}&k\le t<m,\ k\le i\le t,\\
 a_i&k\le t<m,\ i\notin[k,t],\\
 a_i&t\ge m,\ i<k,\\
 a_{i+1}&t\ge m,\ k\le i<m,\\
 1-m&t\ge m,\ i=m,
 \end{cases}
\end{equation}
\begin{equation}\label{eq:older-y}
 y_i(t)=\begin{cases}
 b_{i+1}&t<k,\ i\le t,\\
 b_i&t<k,\ i>t,\\
 b_{i+1}&t\ge k,\ i<k,\\
 a_k-2&t\ge k,\ i=k,\\
 b_i&t\ge k,\ i>k.
 \end{cases}
\end{equation}
For $0\le t<T$ put $u=t-\delta$ and set
\begin{equation}\label{eq:distinct-A}
 A_i(t)=\begin{cases}x_{i+1}(t)+t&r\le u,\ r\le i\le u,\\x_i(t)+t&\text{otherwise},\end{cases}
\end{equation}
\begin{equation}\label{eq:distinct-B}
 B_i(t)=\begin{cases}
 y_i(t)+t&u<0,\\
 y_{i+1}(t)+t&0\le u<r,\ i\le u,\\
 y_i(t)+t&0\le u<r,\ i>u,\\
 y_{i+1}(t)+t&u\ge r,\ i<r,\\
 a_r+t-2&u\ge r,\ i=r,\\
 y_i(t)+t&u\ge r,\ i>r.
 \end{cases}
\end{equation}
Since $u<m$ here, no undefined lower source is read. At $t=T$ use
\begin{align}
 A_i(T)&=\begin{cases}
 a_i+T&i<r,\\
 a_{i+1}+T&r\le i<k-1,\\
 a_{i+2}+T&k-1\le i\le m-2,\\
 \delta+1&i=m-1,\\1&i=m,
 \end{cases}\label{eq:distinct-terminal-A}\\
 B_i(T)&=\begin{cases}
 b_{i+2}+T&i<r,\\a_r+T-2&i=r,\\b_{i+1}+T&r<i<k,\\
 a_k+T-2&i=k,\\b_i+T&i>k.
 \end{cases}\label{eq:distinct-terminal-B}
\end{align}
Empty index ranges are omitted. These formulas include $m=2$, $r=1$, $k=2$, $\delta=1$.

\paragraph{The inequalities used in the verification.}
The hypotheses concern the \emph{original} endpoints. In particular they give
\begin{align*}
 a_i+H&\le b_i &&(1\le i\le m),&
 a_{i+1}+H&\le a_i &&(i<m),\\
 a_i+H&\le b_{i+1} &&(i<k),&
 b_{k+1}+H&\le a_k &&(k<m),\\
 a_i+H&\le b_{i+2} &&(i<r),&
 b_{r+2}+H&\le a_r &&(r+2\le m).
\end{align*}
Also $a_r+H\le b_{r+1}$ because $r<k$, and $a_k+H\le a_r$ because $r<k$. These inequalities follow by combining the definitions of the first crossings with pairwise separation; none applies the separation hypothesis to a shifted endpoint. In particular,
\[
 (a_k-2)-b_{k+1}\ge H-2\quad(k<m),\qquad
 (a_r-2)-b_{r+2}\ge H-2\quad(r+2\le m).
\]
The lower bound after a flip is $H-2$, not $H$. It is positive since $H\ge m+6$. The complementary comparisons are
\[
 b_k-(a_k-2)\ge H+2,\quad
 b_{r+1}-(a_r-2)\ge H+2,\quad
 (a_r-2)-(a_k-2)\ge H.
\]
They keep each flipped upper source between its neighbors, including the case $k=r+1$ in which the two flipped sources are neighbors.

\paragraph{Validity of the candidate rows.}
Write $L_i=x_i(t)+t$ and $V_i=y_i(t)+t$ for the older arrays. For $i<k$, a changed upper source is $b_{i+1}$ and $a_i+H\le b_{i+1}$. At row $k$ after switching, the lower source is at most $a_{k+1}$ if $k<m$, so $a_{k+1}\le b_{k+1}\le a_k-H<a_k-2$. If $k=m$, its new lower endpoint is $t-m+1<a_m+t-2$, since $a_m\ge H$. The remaining rows retain their upper sources and only move lower sources to later indices. Thus $L_i\le V_i$, and the older lower array is nonincreasing with $L_i\le a_i+t$.

The younger operation has three row types. Where it shifts an upper source and $i<r$, we have $L_i=a_i+t$ and $V_{i+1}=b_{i+2}+t$, because $i+1<k$ and $i+1\le t$. Hence $L_i+H\le V_{i+1}$. At its flipped row $r$, the lower source is $L_{r+1}\le a_{r+1}+t\le a_r+t-H<a_r+t-2$. For $i>r$, its lower source is either $L_i$ or $L_{i+1}\le L_i$, while its upper source remains $V_i$. Before the younger switch only the first of these upper-shift cases occurs. This proves row nonemptiness throughout \eqref{eq:distinct-A}--\eqref{eq:distinct-B}. The deletion of entries preserves lower order; the comparisons in the preceding paragraph prove upper order at the two flipped entries, and all other comparisons are inherited or deliberate repetitions.

Every old lower source is at least $H$. The inserted lower value $t-m+1$ is used only for $t\ge m$, so it is positive and is smaller than every surviving old lower source at the same time. The source $b_1$ occurs only at $t=0$; $b_2$ occurs only through $t=\delta$, when $b_2+t=n$. Every other original source is at most $n-\delta-H$, so even at $T=m+\delta$ its translate is at most $n+m-H<n$. A flipped source is two smaller still. These bounds place the arrays in $[1,n]$. At $T$, the new bottom lower values satisfy $1<\delta+1<a_i+T$ for every surviving old lower source, and their upper entries are at least $a_m+T-2>\delta+1$. Thus the terminal arrays are valid as well.

\paragraph{Missing rows.}
The strict original comparisons, with the positive residual gap $H-2$ at flips, leave exactly the displayed repetitions. They give the following complete list:
\begin{center}
\begin{tabular}{lll}
\toprule
Front & Missing ceiling & Missing minimum\\
\midrule
Older & row $t+1$ if $0\le t<k$ & row $t+1$ if $k\le t<m$\\
Younger & row $u+1$ if $0\le u<r$ & row $u+1$ if $r\le u<m$\\
\bottomrule
\end{tabular}
\end{center}
A front that has exited contributes no missing row. The front positions differ by $\delta\ge1$, so a row is never both kinds of defect. For $\delta=1$, two same-kind missing rows can be adjacent. The associated three equal endpoints form a single plateau; the cutoff must skip both missing rows.

\paragraph{Ceiling cutoffs and qualification, including the switches.}
A row with a ceiling has $D_i=U_i=B_i+1$. In a missing-ceiling row, $U_i-1=B_i$, because the upper endpoint repeats the preceding one or belongs to the top plateau at $n$. Its $D_i$ is the next present ceiling below the entire plateau, or zero if there is none.

At the older upper front $h=t+1\le k$, the younger front is above $h$, so
\[
 A_h=a_h+t,\qquad
 D_h=\begin{cases}b_{h+1}+t+1&h<m,\\0&h=m.\end{cases}
\]
For $h<k$, $D_h-A_h\ge H+1$, and the present minimum qualifies for $\Delta C$. For $h=k<m$, $A_h-D_h\ge H-1>0$, so it fails. For $h=k=m$, $A_h\ge H+t>0=D_h$, so it also fails. This accounts for both the ordinary and bottom older switch.

At the younger upper front $h=u+1\le r$, the lower source is $A_h=a_h+t$. The complete cutoff cases are
\[
 D_h=\begin{cases}
 a_k+t-1&t\ge k\text{ and }h=k-1,\\
 b_{h+2}+t+1&\text{otherwise, if }h+2\le m,\\
 0&\text{otherwise, if }h+2>m.
 \end{cases}
\]
In the first case the next present upper source is the already flipped $a_k-2$. In the second case the next source is $b_{h+2}$, including when adjacent ceiling holes require skipping two rows. The third case is the bottom sentinel: it is zero, not $b_{m+1}+t+1$.

If $h<r$, only the second case can occur, since $h+2\le r+1\le k\le m$ and $h\ne k-1$. It gives $D_h-A_h\ge H+1$. At $h=r$ the first case gives $A_h-D_h=a_r-a_k+1\ge H+1$; the second gives $A_h-D_h\ge H-1$; the third gives $A_h>0=D_h$. Thus the younger minimum qualifies exactly for $h<r$ and fails exactly at $h=r$. In particular, when $\delta=1$ and $r=k-1$, the simultaneous switching input is $t=k-1$, $u=r-1$. The older wave has \emph{not} yet switched at that input. Both upper holes skip to the same next ceiling $b_{k+1}+t+1$ if $k<m$, and to zero if $k=m$; both minima fail qualification.

\paragraph{Lower cutoffs and the one-step recurrence.}
Present minima in rows with a ceiling qualify because that ceiling has cutoff $B_i+1>A_i$. Present minima at nonswitching upper holes qualify by the older- and younger-front inequalities just proved. Thus the only failed present minima are the switching rows. Define the set of failed or missing minimum rows at time $t<T$ by
\[
 Q_t=\{t+1:k-1\le t<m\}\ \cup\
     \{u+1:0\le u<m,\ r-1\le u\}.
\]
A set term is omitted when its stated condition fails. For each row $i$, let $s(i)$ be the least index $j\ge i$ outside $Q_t$, if one exists. The minimum cutoff is exactly
\[
 E_i=\begin{cases}A_{s(i)}&s(i)\text{ exists},\\0&\text{otherwise}.
 \end{cases}
\]
For a row outside $Q_t$ this reads $A_i$. A single hole reads $A_{i+1}$, unless it is the bottom row, when it reads zero. For two adjacent holes $i,i+1$, both read $A_{i+2}$, or zero if $i+1=m$. Two separated holes each read the next row. This exhausts the possibilities because $Q_t$ has at most two elements.

For any missing-ceiling row the incomparable remainder in \eqref{eq:row-rule} is $[D_i+1,A_i-1]$ when this interval is nonempty. It joins $[E_i+1,D_i]$. For a row with a ceiling the incomparable part is empty. Therefore in both cases the output endpoints are
\begin{equation}\label{eq:distinct-recurrence}
 A_i(t+1)=E_i+1,\qquad
 B_i(t+1)=\max(D_i,A_i-1).
\end{equation}
Here are the source substitutions in this recurrence. Before an upper switch its minimum qualifies, so the lower endpoint translates. At a switch or existing lower hole, the preceding scan deletes the old lower source and copies the next surviving source. A single hole consequently advances one lower deletion in \eqref{eq:distinct-A}; adjacent holes advance both deletions and produce the same copied value in their two rows. A bottom hole inserts $1$. For $\delta=1$, the two holes at input $t=m-1$ are rows $m-1,m$ and both read zero, producing lower endpoints $1,1$ at time $m$. At that final input only row $m$ is a hole; the next last two lower endpoints are therefore $2,1$, as \eqref{eq:distinct-terminal-A} requires. If $\delta>1$, the older bottom value at the final input is $\delta$ and translates to $\delta+1$ in row $m-1$, while the younger bottom hole inserts $1$ in row $m$. These are the two terminal lower cases; no source $a_{m+1}$ is used.

For upper endpoints, a present ceiling gives $B_i(t+1)=B_i+1$. A nonswitching upper hole has $A_i<D_i$, so its output is the next present upper source plus one, advancing the corresponding deletion in \eqref{eq:distinct-B}. At the older or younger switch, $A_i>D_i$ and the output is $A_i-1$. Relative to translation by $t+1$, its source is respectively $a_k-2$ or $a_r-2$. Adjacent nonswitching upper holes copy the same ceiling, whereas simultaneous switches use their distinct old $A_i$ and the shared lower cutoff. These substitutions give \eqref{eq:distinct-A}--\eqref{eq:distinct-B} at time $t+1<T$, and \eqref{eq:distinct-terminal-A}--\eqref{eq:distinct-terminal-B} at $t+1=T$. Together with the initial arrays they prove the full induction for the actual action.

\paragraph{Rerooting interpretation.}
If the two largest support points belong to distinct edges, their last letters are $BB$: a last $AB$ would be a matched leaf. Their edges are nested. The first balanced-suffix scan identifies $k$ as the mate rank of the last closing. Removing that outer pair and scanning again identifies $r$ as the mate rank of the penultimate closing, giving exactly its defining comparison above. At time $T$, only these two closings have crossed the cut. They become openings at $m+\delta,m$, giving lower endpoints $\delta+1,1$; their former mates become closings, giving upper endpoints $a_k+T-2,a_r+T-2$. All other endpoint labels translate. This is \eqref{eq:distinct-terminal-A}--\eqref{eq:distinct-terminal-B}.

For $\delta\ge m$ and $m\le t\le\delta$, $u\le0$ and the older wave has finished. The displayed arrays reduce to its ordinary first-reroot decoder translated from time $m$. This includes $\delta=m$ and both endpoints of the intermediate range. Thus the macro agrees at every clear cut, proving the remaining cases of \cref{lem:passages}.

\section{Unit-gap collision calculations}\label{app:unit}
This appendix proves the two collision cases of \cref{lem:events}. Unlike the matching passages, these calculations include empty rows. The ceiling and minimum sets used in \eqref{eq:row-rule} are always the genuine sets of the current state.

\subsection{An interior collision}
Let $2\le k\le m-1$, let $b_1=n$, and suppose
\[
 a_k=b_{k+1}+1,\qquad a_i<b_{i+1}\quad(i<k).
\]
Both endpoint lists are initially strict, all other endpoint pairs are separated by at least $H\ge m+6$, and $a_m\ge H$. For $0\le j<k$, use \eqref{eq:ordinary-upper}. At time $j=k$, row $k$ is empty; the other rows have
\[
 A_i(k)=a_i+k,\qquad
 B_i(k)=\begin{cases}b_{i+1}+k&i<k,\\b_i+k&i>k.\end{cases}
\]
For $k+1\le j\le m$, every row is nonempty, with
\begin{align}
 A_i(j)&=\begin{cases}a_{i+1}+j&k\le i\le j-1,\\a_i+j&\text{otherwise},\end{cases}\label{eq:interior-phases}\\
 B_i(j)&=\begin{cases}b_{i+1}+j&i<k,\\a_{k-1}+j-1&i=k,\\b_i+j&i>k.\end{cases}\notag
\end{align}
At $j=m+1$, continue the same upper formula, shift lower sources downward for $k\le i<m$, and put $A_m=1$. This is the terminal formula in \cref{lem:events}.

Before $j=k-1$, the ordinary upper wave applies. At $j=k-1$, the missing ceiling row is $k$ and
\[
 D_k=b_{k+1}+j+1=a_k+j=E_k.
\]
The below-ceiling output interval is empty. Its incomparable interval $[a_k+j+1,b_k+j]$ is contained in the current row, so that output is also empty. All other rows translate.

At time $k$, the empty row has a new ceiling at $(k,a_{k-1}+k)$. It is not dominated by an earlier ceiling, because $a_{k-1}<b_k$, and does not dominate the next-row ceiling, because $a_{k-1}>a_k=b_{k+1}+1$. Every nonempty row contributes its own ceiling and minimum. In the empty row,
\[
 D_k=U_k=a_{k-1}+k,\qquad E_k=a_{k+1}+k.
\]
It therefore refills as $[a_{k+1}+k+1,a_{k-1}+k]$, while the other rows translate. The lower endpoints in rows $k,k+1$ now coincide, and the new physical unit gap is between that upper endpoint and the lower endpoint of row $k-1$.

Subsequently all ceiling rows are present and all present minima qualify. At time $j$ the only missing minimum is row $j$, so its cutoff reads the next lower source, or zero at $j=m$. The lower equality propagates and then inserts $1$, giving \eqref{eq:interior-phases} and the terminal state. Separation verifies every displayed nonempty row, excludes unintended equalities, and keeps the translated old endpoints below $n$ through time $m+1$.

\subsection{A top collision}
Now suppose $m\ge2$, $b_1=n$, $a_1=b_2+1$, all other endpoint separations are at least $H\ge m+6$, and $a_m\ge H$. The first phases are:
\begin{itemize}
\item time $1$: row $1$ is empty; row $i\ge2$ is $[a_i+1,b_i+1]$;
\item time $2$: row $1$ is $[a_2+2,n]$; row $i\ge2$ is $[a_i+2,b_i+2]$.
\end{itemize}
For $3\le j\le m$, the arrays are
\begin{align}
 A_i(j)&=\begin{cases}a_{i+1}+j&i\le j-1,\\a_i+j&i>j-1,\end{cases}\label{eq:top-phases}\\
 B_i(j)&=\begin{cases}b_{i+1}+j&i\le j-2,\\b_i+j&i>j-2.\end{cases}\notag
\end{align}
At time $m+1$ they become
\[
 A_i=a_{i+1}+m+1\ (i<m),\quad A_m=1,\qquad
 B_i=b_{i+1}+m+1\ (i<m),\quad B_m=b_m+m+1.
\]
At time $m+2$, row $m$ is empty and every row $i<m$ is $[a_{i+1}+m+2,b_{i+1}+m+2]$. At time $m+3$ the bottom row refills as $[1,a_m+m+2]$, and all preceding rows translate once more, giving the top-collision terminal formula. When $m=2$ the range in \eqref{eq:top-phases} is empty.

Initially $D_1=b_2+1=a_1=E_1$ and $U_1=n+1$, so row $1$ empties exactly as in the interior collision. At time $1$, however, there is no upset element in the empty top row and therefore no ceiling there. Its cutoffs are
\[
 D_1=b_2+2,\qquad U_1=n+1,\qquad E_1=a_2+1.
\]
The two output intervals join, filling it as $[a_2+2,n]$. At time $2$, the only missing ceiling is row $1$ and the only missing minimum is row $2$. In row $1$, $D_1=b_2+3$ and $E_1=a_2+2$; the incomparable tail is removed by the old interval. In row $2$, the lower cutoff is supplied by $a_3+2$ if $m>2$, otherwise zero. This gives the next phase, including the $m=2$ case.

For $3\le j\le m$, the repeated upper endpoint is in rows $j-2,j-1$, and the repeated lower endpoint is in rows $j-1,j$. Thus the missing ceiling row $j-1$ has lower endpoint $a_j+j<D_{j-1}=b_j+j+1$, propagating the upper equality. The missing minimum in row $j$ reads the next lower endpoint, or zero at the bottom, propagating the lower equality. Every other row translates.

At time $m+1$, only the last two upper endpoints coincide; all lower endpoints are strict. Row $m$ has $D_m=E_m=0$, $U_m=b_m+m+2$, and its current interval is exactly the incomparable range $[1,b_m+m+1]$. It empties. At time $m+2$, this empty bottom row creates a ceiling at $(m,a_m+m+2)$, with $D_m=U_m=a_m+m+2$ and $E_m=0$. It refills as claimed.

The initial exceptional pair disappears at the first empty-row event and returns only at the final refill. All other old endpoints are at most $n-H$, and every shift is at most $m+3<H$. Their original separation and the explicit inserted $1$ guarantee positivity, row validity, and exactly the stated equalities. This finishes the global-rule proof of every phase and of \cref{lem:events}.

\bibliographystyle{alpha}
\bibliography{references}
\end{document}
